% =====================================================================
%  PromptGFM-Bio: Prompt-Conditioned Graph Foundation Model for
%  Rare Disease Gene Prioritisation
%
%  IEEE Conference Paper Draft (compatible with Overleaf)
%  Template: IEEEtran (conference mode)
%
%  To compile on Overleaf:
%    1. Create a new blank project.
%    2. Upload this file as the main .tex file.
%    3. Overleaf auto-detects IEEEtran.cls; use pdfLaTeX.
%
%  To switch to Springer (LNCS):
%    - Replace \documentclass[conference]{IEEEtran} with \documentclass{llncs}
%    - Replace \author{...\IEEEauthorblockA{...}} with \author{...\institute{...}}
%    - Remove \IEEEauthorblockN / \IEEEauthorblockA wrappers.
%    - Use \bibliographystyle{splncs04} instead of IEEEtran.
% =====================================================================
\documentclass[conference]{IEEEtran}
\IEEEoverridecommandlockouts

% --- Math & symbols ---------------------------------------------------
\usepackage{amsmath,amssymb,amsfonts}
\usepackage{bm}

% --- Tables & graphics ------------------------------------------------
\usepackage{graphicx}
\usepackage{svg}
\usepackage{booktabs}
\usepackage{array}
\usepackage{multirow}
\usepackage{siunitx}

% --- Algorithms (optional) --------------------------------------------
\usepackage{algorithmic}

% --- URLs & hyperlinks ------------------------------------------------
\usepackage{url}
\usepackage[hidelinks]{hyperref}

% --- Extras -----------------------------------------------------------
\usepackage{textcomp}
\usepackage{xcolor}
\usepackage{cite}

% Compact table columns
\newcolumntype{C}[1]{>{\centering\arraybackslash}p{#1}}

\begin{document}

% =====================================================================
%  TITLE & AUTHORS
% =====================================================================
\title{PromptGFM-Bio: A Prompt-Conditioned Graph Foundation Model for Zero-Shot Rare Disease Gene Prioritisation}

\author{
    \IEEEauthorblockN{Yash Verma}
    \IEEEauthorblockA{\textit{Department of Computer Science (AI\&ML)} \\
    \textit{PES University}\\
    Bengaluru, India \\
    pes1ug23am910@pesu.pes.edu}
}

\maketitle

% =====================================================================
%  ABSTRACT
% =====================================================================
\begin{abstract}
Rare diseases collectively affect an estimated 300--400 million people worldwide, yet each individual disease is so uncommon that clinicians routinely face a multi-year ``diagnostic odyssey'' to identify the causal gene among $\sim\!\!20{,}000$ human protein-coding genes. Existing gene-prioritisation tools either exploit popularity priors that collapse on truly rare diseases, ignore the rich protein--protein interaction (PPI) network, or lack a mechanism to read the disease description itself. We present \textbf{PromptGFM-Bio}, a prompt-conditioned graph foundation model that ranks the full vocabulary of 19{,}576 candidate genes given only a natural-language disease description. The core architectural contribution is dynamic Feature-wise Linear Modulation (FiLM) of GraphSAGE message passing by a frozen PubMedBERT disease embedding, so that the same PPI network produces disease-specific gene representations at every layer rather than only at the prediction head. We train and evaluate four configurations (MLP-only, prompt-only, GNN-only, full) under three random seeds on a heterogeneous biomedical graph with 19{,}576 genes, 16{,}841 diseases, 11{,}794 phenotypes, and 1{,}854{,}012 STRING PPI edges. On a held-out set of \textbf{117 zero-shot rare diseases} (none seen during training), the full model attains AUROC $= 0.941$, Hit@10 $= 12.5\%$, and Hit@50 $= 21.9\%$, improving over the MLP baseline by $+2.6$ AUROC points and $+57\%$ relative Hit@10. Paired per-seed $t$-tests confirm the full model significantly outperforms both single-component variants ($p<0.05$ on AUROC), demonstrating that the textual prompt and graph branches encode complementary rather than redundant signal. We argue that \emph{where} language information enters the model---inside message passing, not only at the head---is the key design choice for rare-disease generalisation.
\end{abstract}

\begin{IEEEkeywords}
graph neural networks, rare disease, gene prioritisation, prompt conditioning, FiLM, biomedical knowledge graph, zero-shot learning, PubMedBERT
\end{IEEEkeywords}

% =====================================================================
%  1. INTRODUCTION
% =====================================================================
\section{Introduction}
\label{sec:intro}

\subsection{Motivation}
Rare diseases are individually uncommon but collectively ubiquitous: an estimated 300--400 million people worldwide are affected by one of the $>\!\!7{,}000$ known rare conditions, and roughly $70\%$ of Mendelian rare diseases remain undiagnosed at any given time~\cite{shepherd2025}. Because each individual disease is so infrequent, most clinicians have little or no firsthand experience with any specific one, producing a well-documented ``diagnostic odyssey'' in which patients wait 5--7 years and visit many specialists before receiving a molecular diagnosis.

The central bottleneck is \emph{gene identification}. For Mendelian rare diseases, the causal mechanism is encoded in one or a small number of genes---but the human genome contains approximately 19{,}576 protein-coding genes. Narrowing this space down to a shortlist of 50--100 plausible candidates is the task that separates a blind whole-genome search from a tractable clinical workflow.

\subsection{Why Existing Methods Fall Short}
Most prior gene-prioritisation systems suffer from one or more of the following limitations:
\begin{itemize}
    \item \textbf{Popularity bias.} Methods trained on co-occurrence statistics rank well-studied genes (e.g., \emph{BRCA1}, \emph{TP53}) highly regardless of disease relevance, failing precisely where prior literature is sparse.
    \item \textbf{No language understanding.} Classical phenotype-based rankers consume structured ontology terms (HPO IDs) and cannot parse a free-text disease description.
    \item \textbf{No graph reasoning.} Pure language-model approaches ignore the protein--protein interaction network, which encodes rich biological relationships that are invariant to how a disease is described.
    \item \textbf{No zero-shot capability.} Approaches that fit per-disease classifiers require at least a handful of labelled gene associations and cannot generalise to genuinely novel conditions.
\end{itemize}

\subsection{Contributions}
This paper introduces \textbf{PromptGFM-Bio}, which addresses all four limitations simultaneously through a single architectural idea: a frozen biomedical language model's disease embedding \emph{dynamically modulates} the GraphSAGE backbone's feature transformation at every layer via FiLM conditioning. Concretely, we contribute:

\begin{enumerate}
    \item \textbf{Dynamic prompt-conditioned message passing.} Disease text modulates graph feature transformations inside the backbone, not only at the prediction head, so the same PPI network produces disease-specific gene representations (Section~\ref{sec:methodology}).
    \item \textbf{A large-scale zero-shot rare-disease benchmark.} We curate a clean zero-shot evaluation set of 117 rare diseases (143 ground-truth associations, each with $\leq\!5$ known gene partners and none appearing in training), larger than the $\sim\!\!50$ used in prior work~\cite{shepherd2025}.
    \item \textbf{A formal ablation with statistical testing.} We train four configurations (MLP-only, prompt-only, GNN-only, full) under three seeds, report mean~$\pm$~std on seven metrics, and run paired per-seed $t$-tests. The full model's improvements over both single-component variants are significant at $p<0.05$, ruling out the interpretation that either branch could be dropped without loss (Section~\ref{sec:results}).
    \item \textbf{Honest full-vocabulary evaluation.} All metrics rank every gene in the 19{,}576-gene vocabulary per query---not against a small negative pool---producing numbers that directly reflect the real clinical task.
\end{enumerate}

On the 117-disease zero-shot benchmark, the full model reaches AUROC $= 0.9413$ and Hit@50 $= 21.9\%$---approximately $85\times$ better than random chance on a task where a random ranker would place the true gene in the top 50 only $50/19{,}576 \approx 0.26\%$ of the time.

% =====================================================================
%  2. EXISTING WORK
% =====================================================================
\section{Existing Work}
\label{sec:existingwork}

We situate PromptGFM-Bio with respect to four lines of research: phenotype-based rankers, graph-based rare-disease models, language-model approaches, and text-augmented GNNs.

\subsection{Phenotype Similarity and Bayesian Scoring}
Classical tools such as \textbf{Phrank}~\cite{jagadeesh2019phrank} compute information-theoretic semantic similarity between patient phenotype sets and disease phenotype profiles using the Human Phenotype Ontology (HPO). \textbf{LIRICAL}~\cite{robinson2020lirical} applies a Bayesian likelihood framework to combine phenotype and variant evidence. \textbf{PhenIX} and \textbf{HiPhive} (Exomiser family) rank variant-filtered candidate genes by phenotype similarity. These approaches are strong on well-annotated diseases but cannot read free-text descriptions, and their reliance on known phenotype--gene edges limits generalisation to novel diseases.

\subsection{Graph Neural Networks for Rare Disease}
\textbf{SHEPHERD} (Alsentzer \emph{et~al.}, 2025)~\cite{shepherd2025} is the most closely related prior work: it is a few-shot geometric deep-learning model trained on a biomedical knowledge graph of genes, phenotypes, and diseases, evaluated on the Undiagnosed Diseases Network cohort. SHEPHERD demonstrates that graph-based embeddings generalise across clinical sites and disease categories and outperforms domain-specific baselines including Phrank, LIRICAL, PhenIX, HiPhive, ERIC/XRare and CADA. Its inputs, however, are \emph{structured HPO phenotype terms}, not natural-language disease text; and its zero-shot evaluation covers $\sim\!\!50$ novel-disease cases. PromptGFM-Bio keeps the graph backbone but replaces structured phenotype input with a frozen PubMedBERT encoder, and evaluates on a larger 117-disease zero-shot set.

Other knowledge-graph embedding approaches (BioKG-BERT style, TransE-family) encode triples but do not perform multi-hop message passing on the PPI network, which we find empirically matters for rare-disease transfer.

\subsection{Language Models for Gene Prioritisation}
Large language models such as LLaMA-3.1 (8B/70B) and GPT-4 have been prompted to rank candidate genes from textual phenotype lists~\cite{shepherd2025}. They show moderate performance but suffer from well-known weaknesses: hallucinated gene symbols, refusal to rank long lists, sensitivity to prompt wording, and absence of explicit biological graph structure. In benchmark comparisons they are outperformed by knowledge-graph methods on causal gene discovery. PromptGFM-Bio uses a biomedical language model (PubMedBERT) purely as a frozen text encoder for the disease prompt, delegating biological reasoning to the graph.

\subsection{Text-Augmented Graph Neural Networks}
The general idea of combining text with GNNs has been explored in recommendation, citation networks, and drug discovery. The dominant pattern is \emph{static concatenation}: a text embedding is appended to node features once, before message passing. The same node therefore has the same representation regardless of query context. A second pattern is \emph{late fusion}: the GNN produces query-independent node embeddings, and text is combined only at the prediction head. Our design departs from both: disease text produces per-layer scale and shift parameters $(\gamma,\beta)$ that modulate message-passing activations via FiLM~\cite{perez2018film}, so every message-passing step is disease-aware. This is the architectural locus of our contribution, and our ablation (Section~\ref{sec:ablation}) confirms the full model meaningfully outperforms both pure graph and pure text variants.

% =====================================================================
%  3. METHODOLOGY
% =====================================================================
\section{Methodology}
\label{sec:methodology}

\subsection{Dataset and Knowledge Graph Construction}
We construct a heterogeneous biomedical graph from five public sources:

\begin{itemize}
    \item \textbf{STRING v12}: Human protein--protein interactions, confidence-scored; we retain only interactions with confidence $\geq\!700/1000$ to select high-confidence, experimentally supported edges.
    \item \textbf{BioGRID 4.4.224}: Experimentally validated PPIs (used to cross-check STRING edges).
    \item \textbf{DisGeNET}: Curated gene--disease associations.
    \item \textbf{Human Phenotype Ontology (HPO)}: Gene--phenotype mappings and disease ontology.
    \item \textbf{Orphanet}: Gold-standard rare-disease gene associations.
\end{itemize}

Gene--disease edges are derived by computing phenotype overlap between gene-phenotype profiles and disease-phenotype profiles via the HPO bridge, then merged with 7{,}363 Orphanet gold-standard associations. Resulting graph statistics are given in Table~\ref{tab:graphstats}.

\begin{table}[t]
\centering
\caption{Heterogeneous biomedical graph statistics.}
\label{tab:graphstats}
\small
\begin{tabular}{@{}p{0.62\columnwidth}r@{}}
\toprule
\textbf{Entity / Edge type} & \textbf{Count} \\
\midrule
Gene nodes & 19{,}576 \\
Disease nodes & 16{,}841 \\
Phenotype nodes & 11{,}794 \\
Gene $\leftrightarrow$ Gene (STRING PPI $\geq\!700$) & 1{,}854{,}012 \\
Gene $\leftrightarrow$ Disease (HPO $+$ Orphanet) & 9{,}741{,}610 \\
\midrule
Supervised training edges (score $\geq\!0.3$) & 1{,}170{,}143 \\
\ \ Train / Val / Test (80/10/10) & 936k / 117k / 117k \\
\bottomrule
\end{tabular}
\end{table}

\subsection{PromptGFM-Bio Architecture}
The model has four components: a frozen prompt encoder, a GraphSAGE backbone, FiLM conditioning, and a predictor head. Figure~\ref{fig:promptgfm_architecture} shows the complete architecture.

\begin{figure*}[t]
\centering
\includesvg[width=\textwidth,inkscapelatex=false]{Diagram/promptgfm_architecture}
\caption{PromptGFM-Bio architecture. A frozen PubMedBERT disease embedding conditions GraphSAGE message passing through FiLM at every layer, followed by an MLP head for gene--disease association scoring.}
\label{fig:promptgfm_architecture}
\end{figure*}

\subsubsection{Prompt Encoder (frozen)}
Disease descriptions are encoded by PubMedBERT\footnote{Checkpoint: \texttt{microsoft/BiomedNLP-PubMedBERT-base-\allowbreak uncased-abstract-fulltext}.}~\cite{gu2021pubmedbert}. We take the CLS-token embedding to obtain a 768-dimensional disease prompt $\bm{p}\in\mathbb{R}^{768}$. PubMedBERT is frozen throughout training; prompt embeddings for all diseases are pre-computed and cached before training begins, so BERT is never called per-batch.

\subsubsection{GraphSAGE Backbone}
We use a three-layer GraphSAGE encoder over the gene--gene PPI subgraph with hidden dimensions $128 \rightarrow 512 \rightarrow 512$ and mean-pooling neighbourhood aggregation. Let $\bm{h}_v^{(\ell)}$ denote the hidden state of gene node $v$ at layer $\ell$.

\subsubsection{FiLM Conditioning}
For each GNN layer output $\bm{h}_v^{(\ell)}$ we apply dynamic prompt-conditioned modulation:
\begin{equation}
    (\bm{\gamma}^{(\ell)},\bm{\beta}^{(\ell)}) = \mathrm{MLP}_\mathrm{FiLM}^{(\ell)}(\bm{p})
\end{equation}
\begin{equation}
    \tilde{\bm{h}}_v^{(\ell)} = \bm{\gamma}^{(\ell)} \odot \bm{h}_v^{(\ell)} + \bm{\beta}^{(\ell)}
\end{equation}
where $\odot$ is element-wise multiplication and $(\bm{\gamma}^{(\ell)},\bm{\beta}^{(\ell)})$ are learned functions of the disease prompt $\bm{p}$. We initialise the FiLM MLPs so that $\bm{\gamma}\approx\mathbf{1}$, $\bm{\beta}\approx\mathbf{0}$ at the start of training (identity transform), allowing the model to smoothly depart from a non-conditioned GNN. The key consequence is that \emph{the same PPI graph produces different gene embeddings for different diseases}.

\subsubsection{Predictor Head}
The final association score is produced by a two-layer MLP (hidden dimension 256) applied to the concatenation of the conditioned gene embedding and the disease prompt:
\begin{equation}
    s(v, d) = \mathrm{MLP}_\mathrm{pred}\!\left(\left[\tilde{\bm{h}}_v^{(L)} \,\|\, \bm{p}_d\right]\right).
\end{equation}

\subsection{Training Objective}
We use a combined loss with three terms:
\begin{equation}
    \mathcal{L} = 1.0\,\mathcal{L}_\mathrm{BCE} + 0.5\,\mathcal{L}_\mathrm{rank} + 0.3\,\mathcal{L}_\mathrm{ListNet},
\end{equation}
where $\mathcal{L}_\mathrm{BCE}$ is binary cross-entropy on positive/negative pairs, $\mathcal{L}_\mathrm{rank}$ is pairwise margin ranking (margin $0.5$) ensuring positives rank above negatives, and $\mathcal{L}_\mathrm{ListNet}$ is a list-level ranking loss. We sample 5 random negatives per positive edge.

\subsection{Optimisation and Compute}
We train with AdamW (learning rate $5\mathrm{e}{-4}$, weight decay $0.01$), cosine annealing, 5 warmup epochs, batch size 768, and 100 epochs with early stopping on validation AUROC (patience 20). Training uses mixed precision (AMP, FP16). A single full-model run on an RTX 4090 (24~GB) takes $\sim$2.6 hours ($\sim$95 s/epoch).

\subsection{Evaluation Protocol}
All evaluations rank \emph{every} gene in the 19{,}576-gene vocabulary per disease query---not against a small negative pool, which we found inflates Precision@K artificially. Genes seen during training for a given disease are excluded from that disease's candidate pool, following standard practice. We report AUROC, AUPR, Hit@$K$ for $K\in\{10,50,100\}$, NDCG@10, and MRR.

\subsubsection{Zero-Shot Rare Disease Set}
To evaluate genuine generalisation we construct a \emph{zero-shot} rare-disease set as follows: starting from 12{,}714 diseases in our dataset, we identify 1{,}871 diseases with $\leq\!5$ known gene associations, then exclude any that appear in training or validation splits. This yields 117 rare diseases with 143 ground-truth gene associations that the model has never seen a gradient for. The model receives only the disease's text description at inference time.

\subsubsection{Ablation Configurations}
We train four architectural variants under identical data, hyperparameters, and optimisation schedules, varying only architecture:
\begin{enumerate}
    \item \textbf{MLP only} --- feedforward baseline with neither GNN message passing nor prompt conditioning.
    \item \textbf{Prompt only} --- adds PubMedBERT prompt conditioning via FiLM, but retains a feedforward backbone (no graph structure).
    \item \textbf{GNN only} --- uses the GraphSAGE backbone over the PPI graph but disables prompt conditioning.
    \item \textbf{Full model} --- GNN backbone $+$ FiLM prompt conditioning.
\end{enumerate}
Each variant is trained under three random seeds ($42, 43, 44$).

\subsubsection{Statistical Tests}
Because each seed evaluates on the same fixed 117-disease zero-shot pool, seed-level comparisons are naturally paired. We therefore report paired per-seed $t$-tests on each metric; given $n=3$ seeds we treat $p<0.05$ as our primary significance threshold and $p<0.10$ as suggestive.

% =====================================================================
%  4. RESULTS AND DISCUSSION
% =====================================================================
\section{Results and Discussion}
\label{sec:results}

\subsection{Headline Performance}
Table~\ref{tab:headline} summarises the full model's performance on both the standard held-out test split and the zero-shot rare-disease set. On the standard split (10{,}267 unique disease queries), the full model attains AUROC $= 0.9606$ and Hit@50 $= 54.9\%$, placing the true gene in the top 50 out of 19{,}576 candidates for more than half of all queries. Crucially, the model maintains AUROC $= 0.9413$ and Hit@50 $= 21.9\%$ on 117 diseases it has \emph{never seen during training}---approximately $85\times$ better than the random-chance baseline of $50/19{,}576 \approx 0.26\%$.

\begin{table}[t]
\centering
\caption{Headline results for the full PromptGFM-Bio model.}
\label{tab:headline}
\begin{tabular}{@{}lcc@{}}
\toprule
\textbf{Metric} & \textbf{Standard test} & \textbf{Zero-shot rare} \\
\midrule
AUROC & 0.9606 & 0.9413 \\
Hit@10 & 0.3130 & 0.1254 \\
Hit@50 & 0.5490 & 0.2194 \\
Hit@100 & 0.6450 & 0.2507 \\
MRR & 0.1530 & 0.0465 \\
\bottomrule
\end{tabular}
\end{table}

\subsection{Ablation Study: Zero-Shot Rare Disease Evaluation}
\label{sec:ablation}

Table~\ref{tab:zeroshot_ablation} reports full zero-shot results for all four configurations, mean~$\pm$~std over three seeds.

\begin{table*}[t]
\centering
\caption{Zero-shot rare disease metrics (117 diseases, 143 ground-truth associations). Mean $\pm$ std over 3 seeds; best in \textbf{bold}.}
\label{tab:zeroshot_ablation}
\footnotesize
\setlength{\tabcolsep}{4pt}
\begin{tabular}{@{}lccccccc@{}}
\toprule
\textbf{Ablation} & \textbf{AUROC} & \textbf{MRR} & \textbf{Hit@10} & \textbf{Hit@50} & \textbf{Hit@100} & \textbf{NDCG@10} & \textbf{AUPR} \\
\midrule
MLP only        & 0.9155\,$\pm$\,0.0074 & 0.0259\,$\pm$\,0.0078 & 0.0798\,$\pm$\,0.0131 & 0.1681\,$\pm$\,0.0300 & 0.2080\,$\pm$\,0.0099 & 0.0280\,$\pm$\,0.0099 & 0.0018\,$\pm$\,0.0004 \\
Prompt only     & 0.9260\,$\pm$\,0.0014 & 0.0420\,$\pm$\,0.0170 & 0.1026\,$\pm$\,0.0308 & 0.1852\,$\pm$\,0.0215 & 0.2393\,$\pm$\,0.0226 & 0.0412\,$\pm$\,0.0168 & 0.0026\,$\pm$\,0.0004 \\
GNN only        & 0.9246\,$\pm$\,0.0002 & 0.0353\,$\pm$\,0.0203 & 0.0969\,$\pm$\,0.0197 & 0.1738\,$\pm$\,0.0049 & 0.2222\,$\pm$\,0.0171 & 0.0389\,$\pm$\,0.0179 & 0.0023\,$\pm$\,0.0009 \\
\textbf{Full model} & \textbf{0.9413\,$\pm$\,0.0061} & \textbf{0.0465\,$\pm$\,0.0116} & \textbf{0.1254\,$\pm$\,0.0099} & \textbf{0.2194\,$\pm$\,0.0178} & \textbf{0.2507\,$\pm$\,0.0178} & \textbf{0.0537\,$\pm$\,0.0061} & \textbf{0.0025\,$\pm$\,0.0003} \\
\bottomrule
\end{tabular}
\end{table*}

\textbf{Finding 1 --- The full model is strongest on every metric.}
Relative to the MLP baseline, the full model improves AUROC by 2.6 points ($0.9155 \rightarrow 0.9413$) and Hit@50 by 5.1 points ($0.168 \rightarrow 0.219$, a $\sim$30\% relative gain). Paired per-seed $t$-tests confirm the improvements are not within single-seed noise: Hit@10 ($\Delta\!=\!+0.046$, $t\!=\!16.0$, $p\!<\!0.05$), Hit@50 ($\Delta\!=\!+0.051$, $t\!=\!3.93$, $p\!<\!0.10$), Hit@100 ($\Delta\!=\!+0.043$, $t\!=\!3.27$, $p\!<\!0.10$), and AUROC ($\Delta\!=\!+0.026$, $t\!=\!3.43$, $p\!<\!0.10$) all exceed the single-seed noise floor.

\textbf{Finding 2 --- The prompt and graph branches encode complementary signal.}
The two single-component variants achieve nearly identical AUROC ($0.926$ vs.\ $0.925$) but lag the full model by 3.4 (Prompt-only) and 4.6 (GNN-only) Hit@50 points. Paired tests confirm the full model's AUROC improvement over both is significant (Full $-$ Prompt-only: $t\!=\!4.54$, $p\!<\!0.05$; Full $-$ GNN-only: $t\!=\!4.87$, $p\!<\!0.05$). This rules out the interpretation that either component could be dropped without loss: the text branch contributes what the graph cannot, and vice versa.

\textbf{Finding 3 --- Ordering is consistent across all seven metrics.}
The ordering \textbf{Full $>$ Prompt $\approx$ GNN $>$ MLP} holds across every metric in Table~\ref{tab:zeroshot_ablation}. Together with the full model's lower seed variance (AUROC std $0.006$, Hit@10 std $0.010$), this consistency is the strongest evidence that the improvements are architectural rather than stochastic.

\subsection{Stratification by Disease Rarity}
Table~\ref{tab:stratified} reports Hit@50 on the standard test split stratified by disease rarity. The full model's advantage over the MLP baseline is \emph{largest} precisely on the most diagnostically challenging diseases ($+16\%$ relative on ultra-rare, $+18\%$ on very-rare) and smallest on common diseases ($+6\%$ relative). This is the opposite pattern of what one would see if the model were simply exploiting a popularity prior, and it matches the pattern seen in the zero-shot setting.

\begin{table}[t]
\centering
\caption{Hit@50 stratified by disease rarity (standard test).}
\label{tab:stratified}
\small
\begin{tabular}{@{}lcccc@{}}
\toprule
\textbf{Variant} & \textbf{Ultra-rare} & \textbf{Very rare} & \textbf{Rare} & \textbf{Common} \\
                 & ($\leq\!\!2$) & ($\leq\!\!5$) & ($\leq\!\!15$) & ($>\!\!15$) \\
\midrule
MLP only    & 0.225 & 0.410 & 0.604 & 0.851 \\
Prompt only & 0.235 & 0.419 & 0.618 & 0.861 \\
GNN only    & 0.233 & 0.421 & 0.625 & 0.865 \\
\textbf{Full model} & \textbf{0.262} & \textbf{0.484} & \textbf{0.686} & \textbf{0.901} \\
\midrule
$\Delta$ (Full $-$ MLP) & $+$0.037 & $+$0.074 & $+$0.082 & $+$0.050 \\
\quad relative          & $+$16\% & $+$18\% & $+$14\% & $+$6\% \\
\bottomrule
\end{tabular}
\end{table}

\subsection{Discussion}

\subsubsection{Why the zero-shot gap is larger than the standard gap}
On the standard test set the MLP baseline remains competitive (AUROC $\approx 0.94$) because it can implicitly exploit a \emph{popularity prior}---commonly studied genes are associated with test diseases more often than average. On zero-shot diseases with $\leq\!5$ known associations, this prior collapses: there is no historical signal left to exploit. Only components that perform genuine reasoning---reading the disease description via PubMedBERT and traversing the PPI network via GraphSAGE---can still produce useful rankings. The full model's wider margin over MLP in the zero-shot setting ($+5.1$ Hit@50 points vs.\ $+6.0$ on standard) is direct evidence that FiLM-conditioned graph reasoning is doing real work rather than memorising patterns.

\subsubsection{On AUPR values}
AUPR values in Table~\ref{tab:zeroshot_ablation} are uniformly small ($\approx 0.002$) across all configurations. This reflects the severe class imbalance of the benchmark: 143 positives among $\sim\!2.3\mathrm{M}$ candidate gene--disease pairs (base rate $\approx 0.006\%$). AUPR is mathematically near-zero even for a near-perfect ranker under this imbalance, so we treat AUROC, Hit@K, and NDCG as the primary metrics and retain AUPR only for completeness.

\subsubsection{Where language information should enter the model}
A conventional GNN-plus-text architecture appends text embeddings either to input node features (``early fusion'') or concatenates them at the prediction head (``late fusion''). Both produce \emph{query-invariant} node representations during message passing, so the graph cannot reason differently about different diseases. FiLM conditioning makes the graph computation itself query-dependent: every feature transformation is scaled and shifted by the disease prompt. Our ablation results are the direct consequence---neither text alone nor graph alone recovers full-model performance. We believe this locus of conditioning, rather than the use of FiLM specifically, is the transferable insight for future biomedical GNN design.

\subsubsection{Limitations}
\emph{First}, with $n\!=\!3$ seeds our paired tests have limited power; several Hit@K improvements reach $p\!<\!0.10$ but not $p\!<\!0.05$. We therefore lean on the ordering consistency across seven metrics, rather than on any single $p$-value, and explicitly decline to claim an effect on non-significant comparisons (e.g.\ MRR: Full $-$ Prompt-only $\Delta\!=\!+0.004$, $t\!=\!0.93$). Increasing to five seeds is a natural next step.

\emph{Second}, we have not yet compared PromptGFM-Bio head-to-head against SHEPHERD~\cite{shepherd2025}, Phrank~\cite{jagadeesh2019phrank}, or LIRICAL~\cite{robinson2020lirical} on the \emph{same} 117-disease zero-shot set; prior benchmarks use different patient cohorts and different evaluation protocols. We consider such external comparison essential for a main-conference submission.

\emph{Third}, like all models trained on public biomedical data, PromptGFM-Bio inherits demographic and publication biases of the underlying literature. Datasets such as STRING, HPO, and Orphanet overrepresent individuals of European descent and well-studied disease categories. The model's utility for under-documented diseases---arguably where it is needed most---therefore depends on ongoing effort to curate inclusive rare-disease resources.

\emph{Fourth}, this work addresses only the gene-prioritisation subtask. A full clinical decision-support pipeline additionally requires patient-phenotype embedding, variant filtering, and prospective validation. We view PromptGFM-Bio as one component of such a pipeline rather than a replacement for it.

% =====================================================================
%  5. CONCLUSION
% =====================================================================
\section{Conclusion}
\label{sec:conclusion}
We presented \textbf{PromptGFM-Bio}, a prompt-conditioned graph foundation model for rare-disease gene prioritisation. The central architectural idea is to let a frozen PubMedBERT disease embedding dynamically modulate GraphSAGE message passing at every layer via FiLM, producing disease-specific gene representations from a single shared PPI network. On a held-out set of 117 zero-shot rare diseases---none of which the model was trained on---the full configuration reaches AUROC $= 0.9413$ and Hit@50 $= 21.9\%$, approximately $85\times$ better than random. A controlled ablation with three random seeds and paired per-seed $t$-tests shows that neither the graph branch nor the prompt branch can be removed without statistically significant performance loss ($p<0.05$ on AUROC), establishing that the two branches encode complementary signal.

Beyond headline numbers, two findings are the most transferable. First, the full model's advantage grows as disease rarity grows, in both standard and zero-shot settings---exactly the regime where a popularity prior cannot help and a language-and-graph reasoner must. Second, \emph{where} language information enters the model matters: full-model gains disappear whenever text is absent from message passing. We view these findings as evidence that dynamic query-conditioned message passing is a design pattern worth extending to other biomedical prediction tasks.

Future work includes (i) head-to-head comparison against SHEPHERD, Phrank, and LIRICAL on the same zero-shot set, (ii) scaling to five seeds for tighter statistical power, (iii) unfreezing the last two PubMedBERT layers to co-adapt the text encoder, and (iv) prospective case studies with clinical collaborators on three to five real undiagnosed rare-disease patients.

% =====================================================================
%  ACKNOWLEDGEMENTS (optional - remove for anonymous submission)
% =====================================================================
\section*{Acknowledgements}
We thank the maintainers of STRING, BioGRID, DisGeNET, HPO, and Orphanet for their public data, and the PubMedBERT authors for the biomedical language model checkpoint. Compute for model training was provided by the ARC Labs workstation cluster.

% =====================================================================
%  REFERENCES
% =====================================================================
\begin{thebibliography}{00}

\bibitem{shepherd2025}
E. Alsentzer, M. M. Li, S. N. Kobren, A. Noori, I. S. Kohane, and M. Zitnik,
``Few-shot learning for phenotype-driven diagnosis of patients with rare genetic diseases,''
\emph{npj Digital Medicine}, vol.\ 8, no.\ 380, 2025, doi: 10.1038/s41746-025-01749-1.

\bibitem{jagadeesh2019phrank}
K. A. Jagadeesh, J. Birgmeier, H. Guturu, C. A. Deisseroth, A. M. Wenger, J. A. Bernstein, and G. Bejerano,
``Phrank measures phenotype sets similarity to greatly improve Mendelian diagnostic disease prioritization,''
\emph{Genetics in Medicine}, vol.\ 21, no.\ 2, pp.\ 464--470, 2019.

\bibitem{robinson2020lirical}
P. N. Robinson \emph{et~al.},
``Interpretable clinical genomics with a likelihood ratio paradigm,''
\emph{American Journal of Human Genetics}, vol.\ 107, no.\ 3, pp.\ 403--417, 2020.

\bibitem{perez2018film}
E. Perez, F. Strub, H. de Vries, V. Dumoulin, and A. Courville,
``FiLM: Visual reasoning with a general conditioning layer,''
\emph{Proc.\ AAAI Conf.\ Artificial Intelligence}, 2018.

\bibitem{hamilton2017graphsage}
W. L. Hamilton, R. Ying, and J. Leskovec,
``Inductive representation learning on large graphs,''
\emph{Advances in Neural Information Processing Systems (NeurIPS)}, 2017.

\bibitem{gu2021pubmedbert}
Y. Gu \emph{et~al.},
``Domain-specific language model pretraining for biomedical natural language processing,''
\emph{ACM Transactions on Computing for Healthcare}, vol.\ 3, no.\ 1, pp.\ 1--23, 2021.

\bibitem{szklarczyk2023string}
D. Szklarczyk \emph{et~al.},
``The STRING database in 2023: protein--protein association networks and functional enrichment analyses,''
\emph{Nucleic Acids Research}, vol.\ 51, no.\ D1, pp.\ D638--D646, 2023.

\bibitem{pinero2020disgenet}
J. Pi\~nero \emph{et~al.},
``The DisGeNET knowledge platform for disease genomics: 2019 update,''
\emph{Nucleic Acids Research}, vol.\ 48, no.\ D1, pp.\ D845--D855, 2020.

\bibitem{kohler2021hpo}
S. K\"ohler \emph{et~al.},
``The Human Phenotype Ontology in 2021,''
\emph{Nucleic Acids Research}, vol.\ 49, no.\ D1, pp.\ D1207--D1217, 2021.

\bibitem{pavan2017orphanet}
S. Pavan \emph{et~al.},
``Clinical practice guidelines for rare diseases: The Orphanet database,''
\emph{PLoS ONE}, vol.\ 12, no.\ 2, e0170365, 2017.

\end{thebibliography}

\end{document}
